% ---- pg-453 （书页 441） ----

\[
\boldsymbol{\Sigma}=\left[\begin{array}{c:c}\boldsymbol{\Sigma}_{11}&\boldsymbol{\Sigma}_{12}\\ \hdashline \boldsymbol{\Sigma}_{21}&\boldsymbol{\Sigma}_{22}\end{array}\right]=\left[\begin{array}{cc:cc}8&2&3&1\\2&5&-1&3\\ \hdashline 3&-1&6&-2\\1&3&-2&7\end{array}\right]
\]

\noindent (a)\quad 计算典型相关系数 $\rho_1^*,\rho_2^*$.

\noindent (b)\quad 确定典型变量配对 $(U_1,V_1)$ 和 $(U_2,V_2)$.

\noindent (c)\quad 令 $\mathbf{U}=[U_1,U_2]'$ 及 $\mathbf{V}=[V_1,V_2]'$. 由第一准则, 求
\[
E\left(\left[\begin{array}{c}\mathbf{U}\\ \hdashline \mathbf{V}\end{array}\right]\right)\quad\text{和}\quad \mathrm{Cov}\left(\left[\begin{array}{c}\mathbf{U}\\ \hdashline \mathbf{V}\end{array}\right]\right)=\left[\begin{array}{c:c}\boldsymbol{\Sigma}_{UU}&\boldsymbol{\Sigma}_{UV}\\ \hdashline \boldsymbol{\Sigma}_{VU}&\boldsymbol{\Sigma}_{VV}\end{array}\right]
\]
将你得到的结果和结论 10.1 的性质作比较.

\noindent 10.3\quad 令 $\mathbf{Z}^{(1)}=\mathbf{V}_{11}^{-1/2}(\mathbf{X}^{(1)}-\boldsymbol{\mu}^{(1)})$ 和 $\mathbf{Z}^{(2)}=\mathbf{V}_{22}^{-1/2}(\mathbf{X}^{(2)}-\boldsymbol{\mu}^{(2)})$ 为两个标准化变量组, 如果 $\rho_1^*,\rho_2^*,\cdots,\rho_p^*$ 是 $\mathbf{X}^{(1)},\mathbf{X}^{(2)}$ 组的典型相关系数并且 $(U_i,V_i)=(a_i'\mathbf{X}^{(1)},b_i'\mathbf{X}^{(2)}),i=1,2,\cdots,p$, 为相联系的典型变量, 确定 $\mathbf{Z}^{(1)},\mathbf{Z}^{(2)}$ 组的典型相关系数和典型变量, 即用 $\mathbf{X}^{(1)},\mathbf{X}^{(2)}$ 组的典型相关系数和典型变量来表达 $\mathbf{Z}^{(1)},\mathbf{Z}^{(2)}$ 组相应的量.

\noindent 10.4\quad (典型相关系数和变量的交替计算.) 试证明, 如果 $\lambda_i$ 为与特征向量 $\mathbf{e}_i$ 联系的 $\boldsymbol{\Sigma}_{11}^{-1/2}\boldsymbol{\Sigma}_{12}\boldsymbol{\Sigma}_{22}^{-1}\boldsymbol{\Sigma}_{21}\boldsymbol{\Sigma}_{11}^{-1/2}$ 的一个特征值, 则 $\lambda_i$ 还是特征向量为 $\boldsymbol{\Sigma}_{11}^{-1/2}\mathbf{e}_i$ 的 $\boldsymbol{\Sigma}_{11}^{-1}\boldsymbol{\Sigma}_{12}\boldsymbol{\Sigma}_{22}^{-1}\boldsymbol{\Sigma}_{21}$ 的一个特征值.

提示: $|\boldsymbol{\Sigma}_{11}^{-1/2}\boldsymbol{\Sigma}_{12}\boldsymbol{\Sigma}_{22}^{-1}\boldsymbol{\Sigma}_{21}\boldsymbol{\Sigma}_{11}^{-1/2}-\lambda_i\mathbf{I}|=0$ 意味着
\[
0=|\boldsymbol{\Sigma}_{11}^{-1/2}|\,|\boldsymbol{\Sigma}_{11}^{-1/2}\boldsymbol{\Sigma}_{12}\boldsymbol{\Sigma}_{22}^{-1}\boldsymbol{\Sigma}_{21}\boldsymbol{\Sigma}_{11}^{-1/2}-\lambda_i\mathbf{I}|\,|\boldsymbol{\Sigma}_{11}^{1/2}|
\]
\[
=|\boldsymbol{\Sigma}_{11}^{-1}\boldsymbol{\Sigma}_{12}\boldsymbol{\Sigma}_{22}^{-1}\boldsymbol{\Sigma}_{21}-\lambda_i\mathbf{I}|
\]

\noindent 10.5\quad 利用例 10.1 所提供的信息

\noindent (a)\quad 求 $\boldsymbol{\Sigma}_{11}^{-1}\boldsymbol{\Sigma}_{12}\boldsymbol{\Sigma}_{22}^{-1}\boldsymbol{\Sigma}_{21}$ 的特征值并证明这些特征值同 $\boldsymbol{\Sigma}_{11}^{-1/2}\boldsymbol{\Sigma}_{12}\boldsymbol{\Sigma}_{22}^{-1}\boldsymbol{\Sigma}_{21}\boldsymbol{\Sigma}_{11}^{-1/2}$ 的特征值一样.

\noindent (b)\quad 求第二对典型变量 $(U_2,V_2)$, 并由第一准则证明它们的相关系数就是第二典型相关系数 $\rho_2^*=0.03$.

\noindent 10.6\quad 证明将 $\mathbf{X}^{(1)},\mathbf{X}^{(2)}$ 的变量以形式 $\underset{(p\times p)}{\mathbf{C}}\underset{(p\times1)}{\mathbf{X}^{(1)}}$ 及 $\underset{(q\times q)}{\mathbf{D}}\underset{(q\times1)}{\mathbf{X}^{(2)}}$ 做非奇异线性变换后, 典型相关系数是不变的.

提示: 考虑 $\mathrm{Cov}\left(\left[\begin{array}{c}\mathbf{C}\mathbf{X}^{(1)}\\ \hdashline \mathbf{D}\mathbf{X}^{(2)}\end{array}\right]\right)=\left[\begin{array}{c:c}\mathbf{C}\boldsymbol{\Sigma}_{11}\mathbf{C}'&\mathbf{C}\boldsymbol{\Sigma}_{12}\mathbf{D}'\\ \hdashline \mathbf{D}\boldsymbol{\Sigma}_{21}\mathbf{C}'&\mathbf{D}\boldsymbol{\Sigma}_{22}\mathbf{D}'\end{array}\right]$ 及任意线性组合 $a_1'(\mathbf{C}\mathbf{X}^{(1)})=a_1'\mathbf{X}^{(1)}$, 且 $a_1'=a_1'\mathbf{C}$. 类似考虑, $b_1'(\mathbf{D}\mathbf{X}^{(2)})=b_1'\mathbf{X}^{(2)}$, 且 $b_1'=b_1'\mathbf{D}$. 选取 $a_1'=e_1'\boldsymbol{\Sigma}_{11}^{-1/2}\mathbf{C}^{-1}$ 和 $b_1'=f_1'\boldsymbol{\Sigma}_{22}^{-1/2}\mathbf{D}^{-1}$, 使相关系数最大.

\noindent 10.7\quad 令 $\boldsymbol{\rho}_{12}=\begin{bmatrix}\rho&\rho\\\rho&\rho\end{bmatrix}$ 且 $\boldsymbol{\rho}_{11}=\boldsymbol{\rho}_{22}=\begin{bmatrix}1&\rho\\\rho&1\end{bmatrix}$, 对应于等结构相关, 其中 $\mathbf{X}^{(1)}$ 和 $\mathbf{X}^{(2)}$ 每一个都有两个成分.

\noindent (a)\quad 求对应于非零典型相关系数的典型变量.

\noindent (b)\quad 将(a)中结果推广到 $\mathbf{X}^{(1)}$ 有 $p$ 个成分以及 $\mathbf{X}^{(2)}$ 有 $q\ge p$ 个成分的情况.

提示: $\boldsymbol{\rho}_{12}=\rho\mathbf{1}\mathbf{1}'$, $\mathbf{1}$ 是每个元素都为 1 的 $(p\times1)$ 列向量并且 $\mathbf{1}'$ 是每个元素都为 1 的 $(q\times1)$ 行向量. 注意, $\boldsymbol{\rho}_{11}\mathbf{1}=[1+(p-1)\rho]\mathbf{1}$, 故 $\boldsymbol{\rho}_{11}^{-1/2}\mathbf{1}=[1+(p-1)\rho]^{-1/2}\mathbf{1}$.

\noindent 10.8\quad (角测量值的相关系数.) 有某些观测值, 比如风向是由角度表达的. 角 $\theta_2$ 可用一对值 $\mathbf{X}^{(2)}=[\cos(\theta_2),\sin(\theta_2)]'$ 来表示.

\noindent (a)\quad 证明 $b'\mathbf{X}^{(2)}=\sqrt{b_1^2+b_2^2}\cos(\theta_2-\beta)$, 其中 $b_1/\sqrt{b_1^2+b_2^2}=\cos(\beta)$ 且 $b_2/\sqrt{b_1^2+b_2^2}=\sin(\beta)$.
